\section{Approaching the star bound}\label{sec:achieve}

For a fixed placement of the columns, Theorem~\ref{ext:star} asks for a
matrix whose exterior price comes close to the sum of the individual wire
budgets. The rank calculus of Section~\ref{sec:flags} gives two ways to
measure what is lost. Above threshold zero we count column contributions
that become linearly dependent. Below zero we count dimensions that the
columns fail to supply. To turn these losses into a price, we first
recover the exterior value from its threshold ranks.

\subsection{Recovering price from threshold rank}

Let $L=((A_i,y_i,\theta_i))_i$ span $E$, and write $e=\dim E$.
Adjoin a supplier of all of $E$ at the tip with bias $s$, giving value
$\lambda_L(s)$. Write $r_s=r_E(L-s)$ for the rank after subtracting
$s$ from each original bias. By~\eqref{flag:capdef}, the right derivative of
$\lambda_L(s)$ is $e-r_s$. We claim that
\begin{equation}\label{ext:recover}
 \lambda(L)=\int_0^\infty r_s\,ds
             -\int_{-\infty}^0(e-r_s)\,ds .
\end{equation}
Indeed, at high bias the tip can supply all $e$ dimensions, giving
$\lambda_L(s)\ge se$. For an upper bound, transport the ordinary
inputs from the tip. If $a_i=\dim A_i>0$, the degree-$p_i$ input costs
at most $\|y_i\|\log\binom{\ell a_i}{p_i}$. Relative to supply entirely
from the tip, it also changes the bias term by $(\theta_i-s)p_i$.
Maximizing each contribution with the binary entropy bound gives
\[
 0\le\lambda_L(s)-se\le
 \sum_i a_i A(\|y_i\|,\theta_i-s)\longrightarrow0
 \qquad (s\to+\infty).
\]

At low bias the degree-zero tip input gives
$\lambda_L(s)\ge\lambda(L)$. To bound the gain from a tip input of
degree $p$, choose a basis of $E$ from the ordinary suppliers and
assign each basis vector to one containing supplier. In the induced
wedge basis of $E\otimes K^\ell$, the tip input has
$\binom{\ell e}{p}$ coordinates. For each coordinate, wedge its $p$
basis vectors into their assigned ordinary inputs and apply the
coordinate functional at the tip. This term is a restriction of an
ordinary sector. Its reordering sign is constant for the chosen
coordinate and input degrees. The tip contributes bias $ps$, while the
reassigned vectors contribute at least $-pT$ in the ordinary sector,
where $T=\max_i|\theta_i|$. Let $D$ be the diagonal rate of the sites
including the tip, and let $h$ denote binary entropy in nats.
By~\eqref{price:sum}, summing the coordinates costs at most
$D\log\binom{\ell e}{p}$. Hence, for $e>0$,
\[
 0\le\lambda_L(s)-\lambda(L)\le
 e\sup_{0\le u\le1}\{D h(u)+u(s+T)\}
 \longrightarrow0\qquad (s\to-\infty).
\]
Integrating the right derivative of the convex function $\lambda_L$
between these two limits gives~\eqref{ext:recover}. If $e=0$, both
sides vanish.

\subsection{Site flags and the barycenter}

At a site $y$ the threshold spaces $U_{y,-a}$ are nested as $a$
increases. We determine the law of the flag $F(y)$ with strict
thresholds $\{F(y)>a\}=U_{y,-a}$, beginning with $E=K$.
Two suppliers of $K$ at $y$ and the tip, with biases $t$ and $s$,
have value $s+A(\|y\|,t-s)$. Indeed, complementary exterior degrees
$p$ and $\ell-p$ give a perfect pairing of dimension $\binom{\ell}{p}$.
Its price is $\|y\|\log\binom{\ell}{p}$. Taking the entropy limit
and maximizing over $p/\ell$ gives the stated value. Now set $t=-a$.
Subtracting the right derivative at $s=0$ from one gives
\[
 \Pr_{\rm flag}\{F(y)>a\}=(1+\exp(a/\|y\|))^{-1}
 \qquad (\|y\|>0).
\]
Thus $F(y)$ has a logistic law of scale $\|y\|$ and mean zero.
At the tip the value is $\max(-a,s)$, giving the constant flag $F(0)=0$.
The finite second moments of these laws place the flags in the metric
completion of Proposition~\ref{flag:metric}.

The distance between two such flags also has an exterior interpretation.
Take suppliers at $y,z$ with biases $t_1,t_2$. At threshold $s$,
the rank of these suppliers of $K$ is the dimension of the sum of
their threshold spaces. By~\eqref{flag:modular} and~\eqref{flag:pair},
this is the upper tail at $s$ of $\max(F(y)+t_1,F(z)+t_2)$ in the
pair law of $F(y),F(z)$. Equation~\eqref{ext:recover} integrates this
tail to give the expectation of the maximum. On the other hand, the
perfect pairing at $y,z$ gives the value $t_2+A(d(y,z),t_1-t_2)$.
Writing $t=t_1-t_2$ and using $\mathbb E_{\rm flag}F(z)=0$, we obtain
\[
 \mathbb E_{\rm flag}\big(F(y)-F(z)+t\big)_+=A(d(y,z),t).
\]
The right derivative on the left is
$\Pr_{\rm flag}(F(y)-F(z)\ge-t)$. Thus the difference has a logistic
law of scale $d(y,z)$, with a point mass at zero when $y=z$.

The flag metric measures the second moment of this difference, so
$y\mapsto F(y)$ scales every distance by the standard deviation of a
standard logistic variable. Its image is closed by completeness of $Y$
and convex by uniqueness of geodesics in the CAT(0) flag space.
If $y_1,\ldots,y_N$ have barycenter at the tip, the sum of squared
distances to their flags is minimized at $F(0)=0$ within the image.
Projection onto the closed convex image can only decrease these
distances, so zero is also a minimizer in the full flag space.
For any flag $Q$ and $a>0$, comparison of this sum at $aQ$ and at
zero yields
$2a\sum_i\langle Q,F(y_i)\rangle\le Na^2\|Q\|^2$. Letting
$a\downarrow0$ gives
\begin{equation}\label{ext:center}
 \sum_{i=1}^N \langle Q,F(y_i)\rangle\le0
 \qquad\text{for every flag }Q .
\end{equation}
Each summand uses the pair law of $Q$ and the indicated site flag.

\subsection{Controlling lost rank by random insertion}

\begin{proof}[Proof of Theorem~\ref{ext:star}]
Write $b=N/2$. At threshold $s$ each slot $i$ supplies the event
$H_i(s)=\{F(y_i)>s-\theta_i\}$ of dimension
$p_i(s)=\Pr_{\rm flag}(H_i(s))$. For a Vandermonde matrix $C$, let
$r_s(C)$ be the rank of the image of $\bigoplus_i H_i(s)$ under its
column map. Since the sites are unit,
$\int_0^\infty p_i(s)\,ds=A(1,\theta_i)$. Equation~\eqref{ext:recover}
therefore expresses the star deficit exactly as
\[
 \sum_i A(1,\theta_i)-\lambda(C;y,\theta)
   =\int_0^\infty K_+(s)\,ds+\int_{-\infty}^0 K_-(s)\,ds ,
 \qquad
 \begin{aligned}
 K_+(s)&=\sum_i p_i(s)-r_s(C),\\
 K_-(s)&=b-r_s(C).
 \end{aligned}
\]
For $s\ge0$, $K_+$ measures the kernel dimension of the column map
and decreases as its domain shrinks. For $s\le0$, $K_-$ measures the
missing image dimension and decreases as $s$ moves left. Both lie in
$[0,N]$. It suffices to control them at finitely many thresholds for
a single choice of $C$.

\emph{Centering the indicators.}
For $|s|\le M/2$, put
\[
 c_s=\tfrac12-\tfrac{s}{2M},\qquad X_i(s)=\mathbf1_{H_i(s)}-c_s .
\]
The scalar $c_s$ is the fraction of the interval $[-M,M]$ above $s$.
Let $I$ be a uniformly sampled slot. For every fixed finite flag
$Q$, uniformly in this range of $s$,
\begin{equation}\label{ext:meanflag}
 \mathbb E_I\langle Q,X_I(s)\rangle\le\epsilon\|Q\|,
 \qquad \epsilon=q^{-1}+O(N^{-10}).
\end{equation}
To see this, fix one site $y$ and a value $f$ of its flag. Averaging the
centered indicator over a continuous uniform bias on $[-M,M]$ gives
\[
 \min\{1,\max\{0,c_s+f/(2M)\}\}-c_s .
\]
This agrees with $f/(2M)$ on $|f|\le M/2$. By the unit logistic tails
and $M\ge100\log N$, the discrepancy has $L^2$ norm $O(N^{-10})$.
Replacing the continuous bias by the $q$ equally spaced midpoints
changes the average by at most $1/q$ at every $f$. Within the pair law
of $Q,F(y)$, Cauchy--Schwarz bounds the corresponding error in pairing
by $\epsilon\|Q\|$. The remaining terms average to a nonpositive
number by~\eqref{ext:center}, since each site occurs $q$ times.

\emph{A second moment for the insertion sum.}
Choose a uniform permutation $\pi$ of the slots and use it at every
threshold. Fix $s$ for the moment and let $Z_0(s)=0$. Obtain $Z_j(s)$
by inserting $X_{\pi(j)}(s)$ into $Z_{j-1}(s)$ using two-flag
addition. Put $a_j=(\mathbb E_\pi\|Z_j(s)\|^2)^{1/2}$. To handle
sampling without replacement, take a uniform slot $I$ independent of
$\pi$ and exchange it with the slot at position $j$. The new permutation
is uniform, and its $j$th flag $X_I(s)$ is independent of the
\emph{old} partial sum. At most one earlier term changes. Each centered
indicator has norm at most one, and insertion is $1$-Lipschitz in
either argument. Hence the new partial sum lies within distance two of
the old one, even after propagating the change through the intervening
insertions. Lipschitz continuity of the pairing and~\eqref{ext:meanflag}
now give
\[
 \mathbb E_\pi\langle Z_{j-1}(s),X_{\pi(j)}(s)\rangle
    \le 2+\epsilon a_{j-1}.
\]
Expanding the squared norm of each insertion in its two-flag law yields
\[
 a_j^2\le a_{j-1}^2+5+2\epsilon a_{j-1},
 \qquad
 \max_{j\le N} a_j^2=O(N+N^2\epsilon^2).
\]
For the second bound, sum the first inequality up to an index
maximizing $a_j$ and solve the resulting quadratic inequality.

\emph{One matrix for all thresholds.}
For each realized order, apply Lemma~\ref{flag:insertion} with slack one
simultaneously at zero and at $s=\pm M2^{-j}$ for
$1\le j\le\lceil\log_2 N\rceil$. Assign the resulting parameters to
the columns in order $\pi$. This gives one Vandermonde matrix $C=C_\pi$
whose ranks at all selected thresholds satisfy
\[
 r_s(C)\ge\mathbb E_{\rm flag}\min\{b,S_N(s)\}-1 ,
\]
where $S_N(s)$ inserts the uncentered indicators in order $\pi$.
Translation of two-flag addition gives $S_N(s)=Z_N(s)+Nc_s$.
Additivity of the mean at every insertion gives
$\mathbb E_{\rm flag}S_N(s)=\sum_i p_i(s)$. Thus the error above
zero is at most $1+\mathbb E_{\rm flag}(S_N(s)-b)_+$, and the error
below zero is at most $1+\mathbb E_{\rm flag}(b-S_N(s))_+$. At
each selected $s\ne0$ the relevant bound is
\[
 1+\mathbb E_{\rm flag}
   \left(\operatorname{sgn}(s)Z_N(s)-\frac{N|s|}{2M}\right)_+ .
\]
All expectations here use the law of the insertion sum at the single
threshold in question. Averaging over $\pi$ and using
$(x-v)_+\le x^2/(4v)$ for $v>0$, the second-moment bound gives
\[
 \mathbb E_\pi K_\pm(s)\le 1+O\!\left((1+N\epsilon^2)\frac{M}{|s|}\right)
\]
on the respective sides of zero.

Integrate on each side from zero to the nearest grid point using
$K_\pm\le N$, at a cost of at most $M$. On the remaining dyadic
intervals up to $|s|=M/2$, use monotonicity and the bound at the
endpoint nearer zero. Each interval has length equal to that
endpoint's absolute value, so its expected cost is
$O(M(1+N\epsilon^2))$. The interval $M/2\le|s|\le2M$ obeys the
same bound using $|s|=M/2$. For $s>2M$, bound $K_+(s)$ by
$\sum_i p_i(s)$. For $s<-2M$, select any basis of columns of $C$.
The threshold spaces on this basis have images in direct sum, so
$K_-(s)$ is at most the sum of $1-p_i(s)$ on the basis. Since
$|\theta_i|\le M$, the logistic laws give total tail cost $O(Ne^{-M})$.

The expected deficit of the chosen matrices is therefore
$O(M(1+N/q^2)\log N)$. Some order and its matrix attain at most
this average, as required.
\end{proof}
